\section{A preliminary correlation length}\label{sec:preliminary}

The first input is a mixing estimate that is uniform when bonds are
weakened or deleted. Its exponential scale need not be sharp. The strict
inequality in the exponent below is, however, essential to the later
comparison of effective densities. We adapt the pinned-cell argument of
\cite{OpenAI-O4}, retaining its current estimates and modifying the
rotation algebra for spins on $S^2$.

For a bounded local Lipschitz observable $F$ with finite support $A$, set
\[
 \mathcal L(F)=\norm{F}_\infty+\max_{x\in A}\operatorname{Lip}_x(F),
\]
where the individual Lipschitz constant uses round distance on $S^2$.
Distances between supports are measured in the maximum norm, with the
periodic distance on a torus. A finite free subgraph carries no boundary
pins; bonds outside the subgraph are absent. Write $Z_b(n,w)$ for the
homogeneous partition function on the rectangular $n$-by-$w$ torus,
with normalized area measure at every site.

\begin{proposition}[Preliminary mixing and rectangular doubling]
\label{prop:preliminary}
There are constants $b_0,C,c,p>0$ and $\eta>0$ such that, with
\begin{equation}
 X(b)=C\exp\bigl((4\pi-\eta)b\bigr),\qquad b\ge b_0,
 \label{eq:orig-1}
\end{equation}
the following statements hold uniformly for all bond strengths
$0\le\beta_e\le b$, including zero.
\begin{enumerate}[label=\textup{(\roman*)}]
\item On every unpinned rectangular torus whose sides are at least
$CX(b)$, and on every finite free subgraph of the square lattice, local
Lipschitz observables $F,G$ supported on $A,A'$ at distance $d$ satisfy
\begin{equation}
 |\Cov(F,G)|\le
 C(1+b+|A|+|A'|+d)^p\mathcal L(F)\mathcal L(G)
 e^{-cd/X(b)}.
 \label{eq:prelim-mixing}
\end{equation}
In the homogeneous model the periodic infinite-plane limit exists and
is unique.
\item For each fixed aspect bound $\kappa\ge1$, the constants may be
chosen so that, for even $n,w$ with
$\kappa^{-1}\le n/w\le\kappa$ and $\min(n,w)\ge CX(b)$,
\begin{equation}
 0\le4\log Z_b(n,w)-\log Z_b(2n,2w)
 \le C(1+b+n+w)^p e^{-c\min(n,w)/X(b)}.
 \label{eq:prelim-doubling}
\end{equation}
\end{enumerate}
\end{proposition}

The proof has three stages. We first obtain a uniform upper bound on the
response of a pinned square to slowly varying rotations. A one-site
estimate supplies a strict saving in the initial bound. Finally, local
rotation estimates give decay for centered one-site functions, and two
conditional ferromagnetic decompositions extend that decay to arbitrary
local observables.

\subsection{Pinned squares, currents, and rotation identities}

Let $Q(S)=\{0,\ldots,S\}^2$, with arbitrary spins fixed on its boundary.
Orient every bond in a positive coordinate direction. Its coefficient is
$w_e=\beta_e$, except that a bond tangent to the boundary has coefficient
$\beta_e/2$. This allocation makes the actions of adjacent cells add to
the action of their union. For $e=x\to y$, define
\begin{equation}
 s_e=q_x\times q_y,\qquad
 E_e=(q_x\cdot q_y)\Id-\frac12(q_xq_y^T+q_yq_x^T).
 \label{eq:prelim-contact}
\end{equation}
If $F=(F_1,F_2)$ is a smooth $\C^3$-valued test field on the unit square,
$z_e$ denotes the rescaled bond midpoint, and $\mu(e)$ is the bond
direction, put
\begin{align}
 X_Q(F)&=S^{-1}\sum_e w_es_e\cdot F_{\mu(e)}(z_e),
 \label{eq:prelim-current}\\
 \mathcal H_Q(F,G)&=
 S^{-2}\E\sum_e w_e\overline{F_{\mu(e)}(z_e)}\cdot
 E_eG_{\mu(e)}(z_e)-\Cov(X_Q(F),X_Q(G)).
 \label{eq:prelim-stiffness}
\end{align}
The covariance is sesquilinear, with conjugation in the first entry.
The restriction of $\mathcal H_Q$ to constant $3$-by-$2$ matrices is
represented by a real symmetric $6$-by-$6$ matrix $H_Q$. We call this
matrix the stiffness of the pinned square. No lower bound on $H_Q$ is
assumed. For a unit vector $a$,
\[
 a^TE_ea=q_x\cdot q_y-(a\cdot q_x)(a\cdot q_y)
 =q_x^{\perp a}\cdot q_y^{\perp a},
\]
so $|a^TE_ea|\le1$. The half weights give
$\sum_{e\parallel\mu}w_e\le bS^2$, and consequently
\begin{equation}
 H_Q\le b\Id.
 \label{eq:prelim-initial-stiffness}
\end{equation}

Let $R_a(t)$ be ordinary rotation through angle $t$ about the unit axis
$a$. Under $q_x\mapsto R_a(tf_x)q_x$, the action
$\mathcal A=-\sum_e w_eq_x\cdot q_y$ satisfies
\begin{equation}
 D_f\mathcal A=\sum_e w_e(a\cdot s_e)(f_y-f_x),\qquad
 D_f^2\mathcal A=\sum_e w_e(a^TE_ea)(f_y-f_x)^2.
 \label{eq:prelim-score}
\end{equation}
For a general infinitesimal profile $u_x\in\R^3$, direct differentiation
gives
\begin{equation}
 D_us_e=E_e(u_y-u_x)+\frac12(u_y+u_x)\times s_e.
 \label{eq:prelim-bond-derivative}
\end{equation}
Thus, if $u$ vanishes on the pinned boundary and
$(d_Su)_e=S(u_y-u_x)$, integration by parts yields
\begin{equation}
 \E\bigl[\overline{X_Q(d_Su)}X_Q(F)\bigr]
 =S^{-2}\E\sum_e w_e\overline{(d_Su)_e}\cdot E_eF_e
 +\frac12\E X_Q\bigl(F_e\times(\overline{u_y}+\overline{u_x})\bigr).
 \label{eq:prelim-ward}
\end{equation}
Here and below the cross product is extended complex bilinearly. These
are the replacements for \Rref{eq:pre:1-1}, \Rref{eq:pre:1-3}, and
\Rref{eq:pre:1-5}. The current method belongs to the Ward-identity approach to spin
correlations~\cite{AizenmanSimonWard}. The factor $1/2$ in
Equation~\eqref{eq:prelim-ward} determines the variance gain below.

We use an auxiliary parameter $B\ge b$ and impose
\begin{equation}
 \log(2+S)\le16B.
 \label{eq:prelim-size-restriction}
\end{equation}
For every fixed moment order $r<\infty$, the current estimate is
\begin{equation}
 \norm{X_Q(F)}_{L^r}\le
 C_r(1+b)\log(2+S)\norm{F}_{C^1}.
 \label{eq:prelim-current-moment}
\end{equation}
We verify the change in the proof of \Rref{eq:pre:2-2}, because its
uniformity in the pins will be used repeatedly. For a one-axis profile
$f$ that vanishes on the boundary, Equation~\eqref{eq:prelim-score}
and invariance of the product area measure give
\begin{equation}
 \E e^{tD_f\mathcal A}
 \le\exp\left(\frac12bt^2\sum_e(f_y-f_x)^2\right).
 \label{eq:prelim-score-mgf}
\end{equation}
The same Taylor estimate bounds fixed moments of the likelihood ratio
of a profile rotation by $\exp(C_rb\sum_e(f_y-f_x)^2)$.

To extract a transverse current, take a patch of side $d$ with a
proportional buffer, a scalar test $h$ of patch-coordinate $C^1$ norm
$H_{\rm test}$, and a compact score whose potential is affine with slope
$d^{-1}$ in the required spatial direction on the support of $h$.
Choose its color and the rotation axis perpendicular to the desired
current color. In the two opposite rotations use
\[
 f=\arcsin(h/N),\qquad N=C\sqrt{1+b}\,H_{\rm test},
\]
in place of the half angle in \Rref{pre:part-2-2}. Rotation of a bond
current differs from rigid rotation through its midpoint angle by
$O(\sup|df|)$. After multiplication by $N$, the resulting test error is
$O(H_{\rm test}/d)$ per bond. There are $O(d^2)$ bonds and the score
normalization is $d^{-1}$, so its deterministic cost is
$CbH_{\rm test}$. Equation~\eqref{eq:prelim-score-mgf} bounds the moments
of the two rotated scores, including their likelihood ratios, by
$C_r(1+b)H_{\rm test}$. Hence
\[
 \left\|d^{-1}\sum_{e\parallel\mu}w_e(s_e)_k h_e\right\|_{L^r}
 \le C_r(1+b)H_{\rm test}.
\]
A Whitney decomposition of the square has $O(S/d)$ patches at each
boundary-distance scale $d$. Multiplying the patch bound by $d/S$ and
summing the $O(\log(2+S))$ scales proves
Equation~\eqref{eq:prelim-current-moment}; the bounded-width boundary
layer is estimated directly. This is the argument of
\Rref{pre:part-2-3}, with constants independent of the pins and of the
individual strengths.

\subsection{A uniform decrease of stiffness}

The next lemma turns the local identities into a bound at a larger
scale. The enlargement factor is denoted by $\ell$; it is unrelated to
the fixed renormalization factor $L$ used later.

\begin{lemma}[Stiffness iteration]\label{lem:prelim-iteration}
Fix a sufficiently large $D$. Suppose that, on squares of side $R$,
$H_R\le h\Id$ uniformly over all boundary pins and all strengths in
$[0,b]$, where $B^{-D}\le h\le B$ and $B\ge b$ is sufficiently large.
The pinned-cell iteration produces a common side $R_*$ with
\begin{equation}
 H_{R_*}\le B^{-D}\Id,\qquad
 \log(R_*/R)\le4\pi h+C_D\sqrt{B\log B},
 \label{eq:prelim-iteration-conclusion}
\end{equation}
provided Equation~\eqref{eq:prelim-size-restriction} holds throughout
the iteration, including each proposed next square. All stiffness
bounds remain uniform in the original pin and strength class.
\end{lemma}

\begin{proof}
We give the changes to \Rref{pre:cells} through
\Rref{pre:iteration}, including the estimates that determine the
coefficient $4\pi$.

\smallskip\noindent
\emph{The conditional potential and its derivatives.}
Tile a square of side $S=\ell R$ into $\ell^2$ cells and condition on
their complete boundary rings $\omega$.
Their interiors are independent. If $Z_c(\omega)$ is the pinned
partition function in cell $c$, the ring action and conditional current
are
\[
 S_g(\omega)=-\sum_c\log Z_c(\omega),\qquad
 j(F)=\E[X_Q(F)\mid\omega]
     =\ell^{-1}\sum_c\E_cX_c(F).
\]
For the lower-left rescaled anchor $z_c$, put
\[
 N_c(F)=\sup_{c^+}(|F|+\ell^{-1}|\nabla F|),\qquad
 D_c(F)=\ell^{-1}N_c(\nabla F),
\]
where $c^+$ is a fixed enlargement of $c$. Let $P(B)$ denote a
polynomial whose coefficients and degree may depend on fixed moment
orders or requested accuracy, but not on $R,S,\ell$.
Equation~\eqref{eq:prelim-current-moment} gives
\begin{align}
 |\E_cX_c(F)|&\le P(B)N_c(F),\qquad
 |\mathcal H_c(F,G)|\le P(B)N_c(F)N_c(G),
 \label{eq:prelim-cell-bounds}\\
 |\mathcal H_c(F,G)-\overline{F_c}H_cG_c|
 &\le P(B)\{D_c(F)N_c(G)+N_c(F)D_c(G)\}.
 \label{eq:prelim-cell-replacement}
\end{align}
These are \Rref{eq:pre:3-1} and \Rref{eq:pre:3-2}, obtained by
subtracting anchor values in the bilinear form.

For a scalar ring profile $u$, let $H_c^u$ be the cell stiffness with
its pins rotated by $R_a(u)$, and put $O_c(u)=R_a(u(z_c))$, acting on
the color index. Rigid covariance and
Equation~\eqref{eq:prelim-current-moment} give
\begin{equation}
 \|O_c(u)^TH_c^uO_c(u)-H_c\|\le P(B)D_c(u),
 \label{eq:prelim-anchor}
\end{equation}
with the same estimate for the derivative at zero. Indeed, remove the
constant rotation at the anchor, rotate the interior variables by the
remaining profile, and differentiate the contact and covariance in
Equation~\eqref{eq:prelim-stiffness}. The residual score is controlled
by the current estimate; $E_e$ and its first derivative are bounded.
Repeating this calculation at the current pins gives the estimate
along an arbitrary-amplitude path. Each cell law in this calculation
still has the original untwisted interaction. Furthermore,
\begin{equation}
 D_u^2S_g=\ell^{-2}\sum_c
 \mathcal H_c(a\,d_Su,a\,d_Su).
 \label{eq:prelim-ring-hessian}
\end{equation}
These are the required versions of \Rref{eq:pre:3-3} and
\Rref{eq:pre:3-4}.

To exploit the variance in the stiffness, choose $D'>D$ and
\[
 h<v\le2B,\qquad c_0\ge B^{-D'},\qquad
 C_c=v\Id-H_c-c_0P_a\ge B^{-D'}\Id,
\]
where $P_a$ projects onto the two spatial components with color $a$.
Adjoin independent standard Gaussian vectors $n_c\in\R^6$ and
Gaussian vectors $n_{0,c}\in\R^2$ of covariance $c_0\Id$, all independent
of the ring data, and define
\begin{equation}
 J(F)=j(F)+\ell^{-1}\sum_c
       (\sqrt{C_c}\,n_c+a\,n_{0,c})\cdot F_c.
 \label{eq:prelim-augmented-current}
\end{equation}
Total variance gives, for every constant real matrix $A$,
\begin{equation}
 H_Q(A,A)=v|A|^2-\operatorname{Var}J(A).
 \label{eq:prelim-variance-identity}
\end{equation}
This is \Rref{eq:pre:4-3}; it uses no positivity of $H_c$.

Here are the commuting derivatives used to bound the variance. Let
$\psi$ be a nonzero real trigonometric cutoff vanishing on the boundary, and
let $f_i=\psi\varphi_i$, where $\varphi_i$ are real orthonormal Fourier
modes of frequencies at most $4K$, including the constant mode.
For each nonzero pair $\{k,-k\}$, choose one representative $k$ and use
both modes $\sqrt2\cos(k\cdot z)$ and $\sqrt2\sin(k\cdot z)$.
Their respective quadrature modes are
$\sqrt2\sin(k\cdot z)$ and $-\sqrt2\cos(k\cdot z)$.
Frequency sums over both signs therefore agree with sums over these
two normalized real modes. The cutoff has band $O(m)$ and
$K+m\ll\ell$. With $f=(f_i)_i$, define
\[
 G_0=\ell^{-2}\sum_c\nabla f(z_c)\nabla f(z_c)^T,
 \qquad
 p=(c_0G_0)^{-1}\ell^{-1}\sum_c\nabla f(z_c)n_{0,c},
 \qquad \theta=f\cdot p.
\]
Exact quadrature identifies $G_0$ with the continuum gradient Gram
matrix. It is positive definite: a combination of the $f_i$ with zero
gradient is constant and vanishes on the boundary; division by $\psi$
on an open set where it is nonzero then annihilates the Fourier
polynomial, hence every coefficient. The base variables are
\[
 \omega^0=R_a(-\theta)\omega,\quad
 n_{0,c}^0=n_{0,c}-c_0\nabla\theta(z_c)/\ell,\quad
 n_c^0=O_c(\theta)^Tn_c-\sqrt{C_c^0}\,a\nabla\theta(z_c)/\ell,
\]
where $C_c^0$ is evaluated at $\omega^0$. Conditional on these base
variables, the density of $p$ is proportional to $e^{-V(p)}$, with
\begin{equation}
 V(p)=\tfrac12c_0p^TG_0p+S_g(R_a(\theta)\omega^0)
 +\tfrac12\sum_c|n_c^0+\sqrt{C_c^0}\,a\nabla\theta(z_c)/\ell|^2.
 \label{eq:prelim-conditional-potential}
\end{equation}
The successive changes of variables preserve product area measure on
the rings and have constant Gaussian Jacobians. Thus the derivation in
\Rref{pre:gaussian} is unchanged. Differentiation in deterministic
potential directions
$u_i\in\operatorname{span}\{\psi\varphi_j:|k_j|\le4K\}$ at fixed base
variables gives commuting operators $\partial_i$. Their adjoints satisfy
\begin{equation}
 \E\left(\sum_i\partial_i^*Y_i\right)^2
 =\E\sum_{ij}(\partial_jY_i)(\partial_iY_j)
  +\E\sum_{ij}Y_i(\partial_i\partial_jV)Y_j.
 \label{eq:prelim-ibp-array}
\end{equation}
The strictly positive Gaussian quadratic form in
Equation~\eqref{eq:prelim-conditional-potential} justifies integration
by parts, as in \Rref{eq:pre:4-5}.

We now state the band estimates with their accuracy parameters. Fix any
required inverse power $B^{-N}$, and then choose a fixed exponent $J$
sufficiently large. The coefficients and degrees of all polynomial
bounds are fixed before this final enlargement of $J$. Assume $K/\ell\le B^{-J}$, $Bm\le K$, and
Equation~\eqref{eq:prelim-size-restriction}. For fields of band $O(K)$,
\Rref{eq:pre:5-1} and \Rref{eq:pre:5-2} give
\[
 \|N(F)\|_\square\le C\|F\|_2,\quad
 \|D(F)\|_\square\le C(K/\ell)\|F\|_2,\quad
 \|d_Su-\nabla u\|_2\le C(K/S)^2\|\nabla u\|_2,
\]
where $\|a\|_\square^2=\ell^{-2}\sum_c|a_c|^2$.
The Gaussian field $\nabla\theta$ has covariance $c_0^{-1}$ times the
orthogonal projection onto its gradient subspace. Consequently, for
any fixed $n$, outside an event of probability at most $e^{-B^3}$,
\[
 \rho:=\ell^{-1}\|\nabla\theta\|_\infty+
       \ell^{-2}\|\nabla^2\theta\|_\infty\le B^{-n}.
\]
The evaluation bounds for band-limited fields and the grid argument in
\Rref{pre:part-5-1} are scalar spatial estimates and apply without
change. Equations~\eqref{eq:prelim-cell-replacement}--\eqref{eq:prelim-ring-hessian}
then give, simultaneously for all potential directions,
\begin{equation}
 \partial_u^2V\le(v+B^{-N})\|\nabla u\|_2^2
 \label{eq:prelim-band-hessian}
\end{equation}
on this event. The error before choosing $J,n$ is bounded by
$P(B)[K/\ell+(K/S)^2+\rho]\|\nabla u\|_2^2$. Everywhere the same
upper estimate holds with $P(B)$ replacing $v+B^{-N}$, and
\begin{equation}
 \|J(F)\|_{L^r}\le C_rP(B)\ell\|F\|_2.
 \label{eq:prelim-crude-current}
\end{equation}
Thus exceptional events remain negligible even after the fixed powers
of $\ell$ and mode counts used below, since $\log\ell\le16B$.

For real deterministic $F$ of band $O(K)$ perpendicular in color to
$a$, put $F'=-a\times F$.
If at most $CK^2$ real directions in this potential span satisfy
\[
 \left\|\nabla\sum_i d_iu_i\right\|_2\le C|d|,
 \qquad N_c(\nabla u_i)\le C,
\]
the modified nonabelian derivative estimate is
\begin{equation}
 \E\sum_i|\partial_iJ(F)-J(u_iF')|^2
 \le B^{-N}\|F\|_2^2.
 \label{eq:prelim-nonabelian}
\end{equation}
For completeness, if
$r_u(z)=u(z+e_\mu/(2S))+u(z-e_\mu/(2S))-2u(z)$ on direction $\mu$,
the derivative of the physical conditional mean is
\[
 \partial_u j(F)=j(uF')+\tfrac12j(r_uF')
   +\ell^{-2}\sum_c\mathcal H_c(a\,d_Su,F).
\]
This replaces the coefficients $2,1$ in \Rref{eq:pre:5-8} by $1,1/2$.
The Taylor remainder is bounded by
$P(B)\ell KS^{-2}\|\nabla u\|_2\|F\|_2$.
The anchor rotation of the noise differentiates with coefficient one.
Its deterministic drift combines with the last displayed contact to
leave $(H_c+C_c)a=(v\Id-c_0P_a)a$, whose pairing with $F$ is zero.
The remaining scalar linear-functional error is bounded by
$P(B)[K/\ell+(K/S)^2+\rho]\|\nabla u\|_2\|F\|_2$.
Taking its operator norm in $(d_i)$ gives the sum of squared errors
without a mode-count factor. The non-anchor derivative of
$\sqrt{C_c}$ contributes at most
$P(B)K^2\ell^{-2}\|F\|_2^2$ after summing over directions, because the
original noises are independent of the ring data. These are exactly
the estimates \Rref{eq:pre:5-9}--\Rref{eq:pre:5-12}, proving
Equation~\eqref{eq:prelim-nonabelian}.

The other required identity is the uncentered Ward estimate. Define
\[
 \mathcal D(F,G)=v\ell^{-2}\sum_c\overline{F_c}\cdot G_c
                -\E\overline{J(F)}J(G).
\]
For $\zeta=r\xi(z)e^{ik\cdot z}/(i|k|)$ vanishing on the boundary,
with $|r|\le1$, $m\le|k|\le CK$, cutoff amplitude $\xi$ of band
$O(m)$ and bounded fixed-order polynomial derivative norms, and for a
same-frequency exact gradient or bounded constant-matrix test $G$,
assume also that each amplitude's first derivative divided by $|k|$
is bounded, as in \Rref{pre:part-5-4}. Then
\begin{equation}
 |\mathcal D(d_S\zeta,G)|\le P(B)/|k|+B^{-N}.
 \label{eq:prelim-frequency-ward}
\end{equation}
Apply Equation~\eqref{eq:prelim-ward} before replacing the gradient by
its leading Fourier symbol: its contact cancels the raw defect. The
remaining phases cancel in the bracket, whose test has $C^1$ norm
$P(B)/|k|$. Equations~\eqref{eq:prelim-current-moment} and
\eqref{eq:prelim-cell-replacement} finish the estimate. The changed
bracket coefficient affects only an absolute constant.

\smallskip\noindent
\emph{Auxiliary weighted-current control.}
To lower-bound $\operatorname{Var}J(A)$ for a constant matrix $A$, we
will pair $J(A)-\E J(A)$ with the mean-zero divergence
$\sum_i\partial_i^*Y_i$ of a suitable test array $Y=(Y_i)$.
Equation~\eqref{eq:prelim-ibp-array} splits the squared norm of this
divergence into a Hessian contribution and a crossed derivative
contribution. Equation~\eqref{eq:prelim-band-hessian} already controls
the Hessian. Before choosing the main test array, we prove a weighted
current bound that will control its crossed derivative contribution,
with constants independent of the number of frequencies. Set
\[
 M=B^J,\quad K=\ell/B^J,\quad\ell\ge B^{10J},\quad
 m=\lceil B^4\rceil,\quad
 \psi(z)=\prod_{\mu=1}^2[1-\cos^{2m}(\pi z_\mu)],\quad\chi=\psi^2.
\]
For a long step put $K_1=\ell/B^{4J}$ and $K_2=\ell/B^{3J}$; for a
short step put $K_1=8M$ and $K_2=256M$. Frequencies belong to
$2\pi\Z^2$, and
\[
 I_\chi=\int\chi^2\ge1-CB^{-2},\qquad
 T=\sum_{M\le|k|\le K_1}|k|^{-2}
   =\frac1{2\pi}\log(K_1/M)+O(M^{-1}).
\]
Define
\[
 w(k)=\sum_{\substack{H\text{ dyadic}\\M\le H\le K_2}}
 H^{-2}e^{-\sqrt{|k|/H}},\qquad T_w=\sum_k w(k).
\]
The scalar convolution bounds in \Rref{eq:pre:6-3} are
\[
 T_w\asymp1+\log(K_2/M),\quad w*w\le CT_ww,\quad
 (m+|k|)^2w(k)\le C,
\]
and $w(k)\ge c(M+|k|)^{-2}$ for $|k|\le K_2/2$.
Shifts of size $O(m)$ change these weights by a bounded factor; their
tails beyond $K-O(m)$ are smaller than any required inverse power of
$B$, including the polynomial factors in $\ell$.

Fix a unit color $r$ perpendicular to $a$, and put $r'=-a\times r$;
then $r'$ is also unit. Use the annular exact
gradients
$G_p^r=r\,d_S(\psi\widetilde\varphi_p/|k_p|)$,
$2M\le|k_p|\le4M$, where $\widetilde\varphi_p$ is the quadrature
Fourier mode defined above, and directions
$u_i=\sqrt{w(k_i)}\psi\varphi_i$ with $|k_i|\le K$.
Writing $b_p=|k_p|^{-2}$ and
$U^r=\sum_{pi}b_p\E J(G_p^ru_i)^2$, the same-frequency Ward bound gives
$\sum_pb_p\E J(G_p^r)^2\le Cv$. In the integration-by-parts proof of
\Rref{eq:pre:6-5}, Equation~\eqref{eq:prelim-nonabelian} supplies signal
$U^r$ rather than $2U^r$. The Hessian contribution is at most $CvU^r$,
and convolution bounds the squared differentiated arrays by
$CT_wU^{r'}$, up to arbitrary inverse powers of $B$.
Thus, for $U=\max(U^r,U^{r'})$,
\[
 U^2\le Cv(v+T_w)U+CvB^{-N_*},\qquad
 U^r+U^{r'}\le Cv(v+T_w),
\]
where $N_*>100(1+D+D')$ is fixed first. The derivative errors cost
$B^{-N}T_w$, not $B^{-N}K^2$. Taking common positive lower weights,
then pairing opposite frequencies and coordinate reflections, proves
\begin{equation}
 \sum_{|n|\le4K_1}(M+|n|)^{-2}
 \sum_{\mu=1}^2\E|J(r e_\mu\psi^2e^{in\cdot z})|^2
 \le Cv(v+T_w),
 \label{eq:prelim-weighted-current}
\end{equation}
as in \Rref{eq:pre:6-6} and \Rref{eq:pre:6-7}.

\smallskip\noindent
\emph{The main divergence test and its variance gain.}
For a constant real $3$-by-$2$ matrix $A$ with $|A|=1$, choose a unit
$a$ normal to both columns, and put $A'=-a\times A$. On the main
annulus take
\[
 u_i=\chi\varphi_i/|k_i|,\quad
 G_i^{A'}=d_S\bigl(\chi(A'\widehat k_i)
                        \widetilde\varphi_i/|k_i|\bigr),\quad
 Y_i=J(G_i^{A'})/|k_i|.
\]
The band separation ensures that these $u_i$ belong to the potential
span. Their gradient Gram operator satisfies
\[
 \left\|\nabla\sum_i d_i u_i\right\|_2^2
 \le(1+Cm/M)^2\sum_i|d_i|^2,
\]
because $\|\chi\|_\infty\le1$, $\|\nabla\chi\|_\infty\le Cm$, and
all main-annulus frequencies are at least $M$.
In each current pairing, first combine the normalized cosine and sine
modes into the equivalent Hermitian pairing of the two complex
frequencies. Apply Equation~\eqref{eq:prelim-frequency-ward} to the exact
discrete gradient before making the leading-symbol replacement in the
remaining contact pairing. Exact quadrature and
$\sum|k|^{-2}\widehat k\widehat k^T=(T/2)\Id$ then give the first two
lines below. The Gram bound and
Equation~\eqref{eq:prelim-band-hessian}, with the negligible exceptional
event estimated by Equation~\eqref{eq:prelim-crude-current}, give the
third:
\begin{align*}
 \sum_i\E Y_i^2&=vI_\chi T/2+O(B^{-N_*}),\\
 \sum_i\E(\partial_iJ(A))Y_i&=vI_\chi T/2+O(B^{-N_*}),\\
 \E Y^TV''Y&\le v^2I_\chi T/2+O(B^{-N_*}).
\end{align*}
Only the second line changes from \Rref{eq:pre:7-1}.
The crossed-array bound remains
\[
 \left|\sum_{ij}\E(\partial_jY_i)(\partial_iY_j)\right|
 \le C\{vT^2+v(v+T_w)\}.
\]
Indeed, the two differentiated arrays each acquire the new constant
one. Their complex pairing is still the sesquilinear pairing at total
frequency $n=k+l$. The convolution
$\sum_{k+l=n}|k|^{-2}|l|^{-2}$ is bounded by
$CT(M+|n|)^{-2}$. The longitudinal part is bounded by
Equation~\eqref{eq:prelim-frequency-ward}; for the transverse part,
\[
 |\sin\angle(k,n)|\,|\sin\angle(l,n)|
 \le\min\{1,\min(|k|,|l|)/|n|\}
\]
removes the convolution logarithm. Equation~\eqref{eq:prelim-weighted-current}
then gives the second term. The cutoff replacements are those of
\Rref{pre:part-7-1}: terms divisible by $\psi^2$ use the weighted bound;
terms not so divisible have an $S^{-2}$ Taylor factor and use
Equation~\eqref{eq:prelim-crude-current}, with $\ell/S^2\le1/\ell$.
No regularity of the random ring data is required.

Pair the mean-zero divergence in Equation~\eqref{eq:prelim-ibp-array}
with $J(A)-\E J(A)$. Cauchy--Schwarz gives the modified form of
\Rref{eq:pre:7-8}:
\begin{equation}
 \operatorname{Var}J(A)\ge
 \frac{(vI_\chi T/2+O(B^{-N_*}))^2}
 {v^2I_\chi T/2+C\{vT^2+v(v+T_w)\}+O(B^{-N_*})}.
 \label{eq:prelim-variance-gain}
\end{equation}
For a long step, $v\ge h\ge1$ and $\log\ell\le h$, this yields
\begin{equation}
 \operatorname{Var}J(A)\ge\frac{\log\ell}{4\pi}
 -C_J\left(\log B+\frac{(\log\ell)^2}{v}\right).
 \label{eq:prelim-long-gain}
\end{equation}
For a short step $T,T_w$ stay between positive absolute constants, so
\begin{equation}
 \operatorname{Var}J(A)\ge c_{\rm sh}\min(v,1).
 \label{eq:prelim-short-gain}
\end{equation}
The constant $c_{\rm sh}>0$ is chosen before the final increase of $J$
and $B$; those increases suppress errors, while all leading constants
in the weighted and crossed estimates are absolute.

\smallskip\noindent
\emph{Summing the gains.}
For $h\ge A_0\log B$, take
$v=h+1$, $c_0=1/2$, and
$\ell=\lceil\exp\sqrt{h\log B}\rceil$, where $A_0$ is fixed large
enough for the band conditions and the error in
Equation~\eqref{eq:prelim-long-gain}. Equations~\eqref{eq:prelim-variance-identity}
and \eqref{eq:prelim-long-gain} decrease $h$ by $\delta h$ with
\[
 \delta h\ge\frac{\log\ell}{4\pi}-C\log B
 \ge c\sqrt{h\log B}.
\]
There are at most $C\sqrt{B/\log B}$ such steps. Summing
$\log\ell\le4\pi\delta h+C\log B$ bounds their total logarithmic
length by $4\pi h+C_D\sqrt{B\log B}$. Below $C\log B$, choose a
fixed $0<\delta<c_{\rm sh}/4$ and take
\[
 v=h+\delta\min(h,1),\quad c_0=\tfrac12\delta\min(h,1),\quad
 \ell=\lceil B^{11J}\rceil.
\]
Equation~\eqref{eq:prelim-short-gain} decreases $h$ by at least
$c\min(h,1)$. The $O_D(\log B)$ remaining steps cost
$O_D((\log B)^2)$ in logarithmic length and reach $B^{-D}$.
This proves Equation~\eqref{eq:prelim-iteration-conclusion}. At every
stage the only conditioned cell laws are the original nearest-neighbor
laws with new boundary pins, so the uniformity class is preserved.
\end{proof}

\subsection{A strict saving on the first square}

Starting Lemma~\ref{lem:prelim-iteration} with
Equation~\eqref{eq:prelim-initial-stiffness} would give logarithmic
length $4\pi b+o(b)$. The next argument saves a fixed multiple of $b$
before the iteration begins.

\begin{lemma}[Initial reduction]\label{lem:prelim-initial-reduction}
There is a fixed sufficiently small $u>0$ such that, for all sufficiently
large $b$, the square of side $S=\lfloor e^{2\pi bu}\rfloor$ satisfies
\[
 H_S\le b(1-3u/4)\Id
\]
uniformly in all its boundary pins and all strengths $0\le\beta_e\le b$.
\end{lemma}

\begin{proof}
At a site at distance at least $S/b$ from the boundary, write its
one-site density relative to normalized area as $e^{-F}$.
For each fixed $u>0$,
\begin{equation}
 \nabla^2F\le(1+o_b(1))u^{-1}g_{S^2},
 \label{eq:orig-2}
\end{equation}
uniformly in the pins, the strengths, and the point of $S^2$.
To prove this, follow any unit-speed great circle through the specified
spin and rotate the other spins about its axis with a radial cutoff.
The cutoff equals one within radius $2$, vanishes at radius $S/(2b)$,
and is linear in the logarithm of the radius between them. The second
derivative of the negative logarithm of the marginal integral is the
expected action second derivative minus a variance. By
Equation~\eqref{eq:prelim-score}, its upper bound is $b$ times the
cutoff Dirichlet energy. Taylor expansion on each bond gives
$\log|x+e_\mu|-\log|x|=e_\mu\cdot\nabla\log|x|+O(|x|^{-2})$.
The summed square error and the integral-comparison error are $O(1)$
before division by the square of the logarithmic radius. Hence the
cutoff energy is
\[
 \frac{2\pi+o_b(1)}{\log S}
 =\frac{1+o_b(1)}{bu},
\]
which proves Equation~\eqref{eq:orig-2}. No positive lower bound on any
coupling entered this calculation.

Fix a unit vector $v$ and put $z=v\cdot q$, $Y=\nabla z$, and
$M=\E z^2$. On $S^2$,
$|Y|^2=1-z^2$, $\operatorname{div}Y=-2z$, and
$\nabla_YY=-zY$. Integration by parts against $e^{-F}$ gives
\begin{equation}
 \E(YF)^2=\E\{\nabla^2F(Y,Y)-3zYF\},\qquad
 \E zYF=1-3M.
 \label{eq:prelim-fisher-identities}
\end{equation}
Writing $a=(1+o_b(1))/u$, the first identity gives
$\E(YF)^2\le a+6$. Cauchy--Schwarz therefore yields
$(1-3M)^2\le M(a+6)$. If $M\le2u$, then
\[
 M\ge\frac{(1-6u)^2}{a+6}
   =u(1-O(u)-o_b(1));
\]
if $M>2u$, the same lower bound is immediate. First choosing $u$ small
and then $b$ large makes $M\ge(7/8)u$ at every interior site.

For an edge with both endpoints in this interior and any unit $v$,
Cauchy--Schwarz gives
\[
 \E[v^TE_ev]
 =\E(q_x^{\perp v}\cdot q_y^{\perp v})
 \le\sqrt{(1-\E(v\cdot q_x)^2)(1-\E(v\cdot q_y)^2)}
 \le1-7u/8.
\]
The excluded boundary strips contain an $O(1/b)$ fraction of the
bonds. Since $0\le w_e\le b$, their contact contribution is bounded
directly. Summing the contacts separately in the two spatial directions
and discarding the nonpositive covariance in
Equation~\eqref{eq:prelim-stiffness} gives
$H_S\le b(1-3u/4)\Id$ for sufficiently large $b$.
\end{proof}

Take $B=b$ and apply Lemma~\ref{lem:prelim-iteration} after this first
square. The total logarithmic side is at most
\begin{equation}
 2\pi bu+4\pi b(1-3u/4)+C_D\sqrt{b\log b}
 =(4\pi-\pi u)b+o(b).
 \label{eq:prelim-total-length}
\end{equation}
The size premise is valid before each application: accumulated
logarithmic length is at most $(4\pi-\pi u)b+o(b)$, and a proposed next
step adds at most $b$. Since $4\pi+1<16$, the premise holds for all
large $b$. Any fixed polynomial enlargement can be absorbed by taking,
for example, $\eta=\pi u/2$ and increasing the threshold for $b$.
Thus there is a common side $R_*$ with
$H_{R_*}\le b^{-D}\Id$ and every later polynomial enlargement bounded
by the length in Equation~\eqref{eq:orig-1}.

\subsection{From small stiffness to local decorrelation}

We next prove that the small-stiffness square controls rotations under
arbitrary conditioning outside a collar. This supplies the centered
one-site bound needed for the final ferromagnetic decomposition.
Choose coarse cells with sides comparable to
$s=R_*\lceil b^{D+C_0}\rceil$. Rectangles with sides in $[s,2s]$
permit tilings of all sufficiently large rectangular tori. Profiles
are supported on a fixed number of cells, vanish at the outer boundary
of a fixed collar, and have uniformly bounded cellwise $C^2$
extensions. All pins lie outside this collar.

For the effective ring action $F_{\rm ring}$, subdivision into
$R_*$-squares and leftover strips gives
\begin{equation}
 D_f^2F_{\rm ring}\le Cb^{-D}
 \label{eq:prelim-profile-hessian}
\end{equation}
for each bounded fixed-axis profile. Indeed, putting
$M_*=\lceil b^{D+C_0}\rceil$, each complete small square costs
$CM_*^{-2}[b^{-D}+P(b)/M_*]$ by its stiffness bound and
Equation~\eqref{eq:prelim-cell-replacement}. The strips have total
area $O(sR_*)$ and cost at most $Cb/M_*$ by the contact bound.
Choose $C_0$ above the degree of the fixed polynomial in the cell
replacement estimate. This is \Rref{eq:pre:9-1}.

Integration by parts gives
$\E(D_fF_{\rm ring})^2\le Cb^{-D}$. The square-root density relative
to product area therefore has rotation derivative of squared
$L^2$ norm at most $Cb^{-D}$. Rotation pullbacks are unitary, so
integration along a fixed-axis path and telescoping finitely many
factors give
\begin{equation}
 \E\left(e^{-[F_{\rm ring}(\phi\omega)-F_{\rm ring}(\omega)]/2}-1\right)^2
 \le Cb^{-D}.
 \label{eq:prelim-profile-displacement}
\end{equation}
For a fixed-dimensional compact family of smooth profiles with bounded
parameter derivatives and a bounded smooth homotopy to identity,
\Rref{pre:part-9-1} upgrades this to
\begin{equation}
 \Pr\!\left(\sup_\lambda
 |F_{\rm ring}(\phi_\lambda\omega)-F_{\rm ring}(\omega)|>\epsilon\right)
 \le\zeta(b),\qquad\zeta(b)\longrightarrow0.
 \label{eq:prelim-profile-uniform}
\end{equation}
The input is Equation~\eqref{eq:prelim-score-mgf}: compact Lie-algebra
scores have Gaussian tails with variance bound $Cb$, and rotated
scores have fixed moments bounded polynomially in $b$. For parameter
dimension $d_0$, Morrey's estimate and a mesh of spacing $b^{-6}$
therefore reduce the supremum to $O(b^{6d_0})$ uses of
Equation~\eqref{eq:prelim-profile-displacement}. Taking
$D>10d_0+10$ makes the error tend to zero. These estimates are uniform
under conditioning outside the collar.

To obtain a finite-range orientation law, adjoin independent Haar
variables $g_X\in S^3$ at free coarse vertices, with identity values at
pinned vertices. The group acts on $S^2$ by the covering homomorphism
$\rho:S^3\to SO(3)$. On a coarse edge $X\to Y$, interpolate by
\[
 \Phi(g;t)=g_X\exp\{\vartheta(t)\ell_0(g_X^{-1}g_Y)\},
\]
where $\vartheta$ is fixed smooth and constant near the endpoints, and
$\ell_0$ is a measurable logarithm of norm at most $\pi$. Write the
retained ring variables as $\omega=\rho(\Phi(g))\widetilde\omega$.
For fixed $g$ this preserves product area measure, and the conditional
law of $g$ has finite-range cell potentials. Common left multiplication
satisfies $\Phi(kg)=k\Phi(g)$; each unpinned cell potential is therefore
invariant under this multiplication.

The profile construction of \Rref{pre:part-10-1} takes place in $S^3$
itself. Fix a coarse vertex and compare two trial assignments differing
only at that vertex, relative to the actual group data. Their relative
profiles agree on the perimeter of the incident cells, which we call
the star. Enlarge the parameter set by treating each vertex quaternion
and each edge logarithm of norm at most $\pi$ as independent variables.
The endpoint equations and agreement on unaffected edges define a
compact subset of valid assignments, containing every chosen measurable
logarithm branch. At each valid assignment, the finitely many relative
edge arcs together with the identity omit some point of $S^3$.
Their distance from this point stays positive on a parameter
neighborhood. In the resulting stereographic chart, correct endpoints
smoothly off the valid subset, with zero correction on that subset.
The interpolation of \Rref{pre:part-10-1} then gives cell extensions
and homotopies to identity, chosen to agree outside the star for the
two trial assignments. A finite cover by such parameter neighborhoods
supplies uniformly bounded smooth families. The supremum over these
families controls every measurable branch, without differentiating
the branch selection.

Compose these $S^3$-valued families with the fixed smooth action
$(g,q)\mapsto\rho(g)q$ on $S^2$. Compactness bounds its required
derivatives, and composition preserves the homotopies and agreement
outside the star. Thus Equation~\eqref{eq:prelim-profile-uniform} applies
to the whole star test: call a coarse vertex bad when varying its group value,
for some neighboring values, changes its incident-cell action by more
than $\epsilon$. A finite coloring of overlapping collars gives, for
each finite vertex set $E$,
\[
 \Pr(E\text{ consists of bad vertices})\le\zeta(b)^{|E|/C}.
\]
At a good vertex each conditional orientation density minorizes Haar
measure by $e^{-\epsilon}$. The disagreement exploration in
\Rref{pre:part-10-2}, with failure probability
$q_0=1-e^{-\epsilon}$, consequently assigns a fixed path of $r$ vertices
probability at most $(q_0+\zeta(b)^{1/C})^r$. Choose $\epsilon$ small
and then $b$ large enough to beat the bounded-degree path count.
This proves exponential decay from the pinned boundary on scale $s$.
The exploration samples original conditional marginals at every reveal,
so its adaptive order preserves both marginal laws.

Averaging the comparison over common left rotations, and using the
transitivity of Haar measure on $S^2$, now gives
\begin{equation}
 |\E O(q_x)O'(q_y)|\le
 C\norm{O}_\infty\norm{O'}_\infty e^{-cd(x,y)/s}
 \label{eq:prelim-one-site}
\end{equation}
for every bounded Haar-centered one-site function $O$ and every bounded
one-site $O'$. To see the passage explicitly, condition on the boundary
of a coarse region around $x$ of radius one third of $d(x,y)$, use the
rotation comparison inside it, and then multiply by $O'(q_y)$.
Small distances use the trivial bound. This is the $S^2$ version of
\Rref{eq:pre:10-4}; it applies also to merely measurable coordinates.
Finite free subgraphs are covered by padding with zero couplings.

\subsection{Neutral observables and completion of the proposition}

The preceding estimate does not require a local observable to be
smooth, but it concerns a centered one-site function. To control all
local observables, write
\[
 q_i=(\cos\alpha_i\,z_i,\sin\alpha_i\,w_i),\qquad
 0\le\alpha_i\le\pi/2,\quad z_i\in S^1,\quad w_i\in\{-1,1\}.
\]
Normalized area measure is the product of uniform circle measure, a
fair sign, and $\cos\alpha\,\dd\alpha$. In particular the Haar mean of
$\alpha$ is $\pi/2-1$. Conditional on $\alpha$, the $z$ and $w$
variables form independent zero-field ferromagnetic XY and Ising
systems with couplings
\[
 J_e^1=\beta_e\cos\alpha_x\cos\alpha_y,
 \qquad J_e^2=\beta_e\sin\alpha_x\sin\alpha_y.
\]
Their partition functions multiply to give the density of $\alpha$
relative to its product one-site measure.

The correlation inequalities used here are the plane-rotator form of
Ginibre's inequality~\cite{Ginibre1970} and the ferromagnetic Ising
inequalities. The XY and Ising energy means and energy covariances are nonnegative,
and their two-point functions are nondecreasing in every coupling;
these inequalities, including zero couplings, are proved in
\Rref{pre:part-11-1}. In each sector all first derivatives of couplings
with respect to the angles have the same sign, and mixed endpoint
derivatives of a single coupling are nonnegative. The mixed
logarithmic derivatives of the angle density are therefore
nonnegative. Its log lattice condition gives association by the
Fortuin--Kasteleyn--Ginibre argument~\cite{FortuinKasteleynGinibre1971}. We use the
associated-variable inequality
\begin{equation}
 |\Cov(f,g)|\le\sum_{ij}
 \operatorname{Lip}_i(f)\operatorname{Lip}_j(g)
 \Cov(\alpha_i,\alpha_j),
 \label{eq:prelim-association}
\end{equation}
proved in \Rref{eq:pre:11-1} by comparing $f,g$ with the increasing
linear functions having their individual Lipschitz constants.

Here is why the quantitative localization in
\Rref{eq:pre:11-3} remains valid with one sign sector. Splitting the XY
angle into an angle in $[0,\pi/2]$ and two signs produces two independent
conditional ferromagnetic Ising systems. The additional sector $w$
is a third such system and has no further angle variable.
Its sign covariance is controlled directly by the associated-variable
inequality, with sign dependence bounded by a constant times the
supremum norm. For the two XY signs, the truncation and angle
localization in \Rref{pre:part-11-2} give the fractional two-point
bound with exponent $1/3$. For the third sign correlation
$0\le c^2_{ij}\le1$, one has $c^2_{ij}\le(c^2_{ij})^{1/3}$.
Thus the same sum of conditional two-point functions to the power
$1/3$, with polynomial support and boundary factors, controls arbitrary
conditional covariances. Deleting conditional Ising bonds gives the
same boundary sums, by coupling monotonicity.

The normalized circle components, $w=\operatorname{sign}(q_3)$, and
$\alpha-(\pi/2-1)$ are bounded Haar-centered functions of one spin.
Equation~\eqref{eq:prelim-one-site} consequently implies, for the full
conditional-sector two-point functions $c^1_{ij},c^2_{ij}$,
\begin{equation}
 0\le\E_\alpha c^t_{ij}\le Ce^{-cd(i,j)/s}\quad(t=1,2),
 \qquad |\Cov(\alpha_i,\alpha_j)|\le Ce^{-cd(i,j)/s}.
 \label{eq:prelim-sector-correlation}
\end{equation}
This is the replacement for \Rref{eq:pre:11-4}.

Let $F,G$ have supports $A,A'$ at distance $d$, and take neighborhoods
of radius $\lfloor d/10\rfloor$. For the conditional covariance given
$\alpha$, apply the preceding arbitrary-observable bound and average.
Jensen's inequality and Equation~\eqref{eq:prelim-sector-correlation}
give $\E_\alpha(c^t_{ij})^{1/3}\le(\E_\alpha c^t_{ij})^{1/3}$,
so this contribution is a polynomial times $e^{-c'd/s}$.

For the covariance of the conditional means, delete the conditional
XY and Ising bonds crossing the chosen neighborhoods, keeping the
original full angle marginal fixed. Each cut mean depends only on its
local angle set. Differentiation with respect to one angle differentiates
at most its incident couplings, each with derivative bounded by $b$;
edge observables have absolute value at most one. Hence the individual
angle Lipschitz constants are at most
$C(1+b)\mathcal L(F)$ and $C(1+b)\mathcal L(G)$.
Equations~\eqref{eq:prelim-association} and
\eqref{eq:prelim-sector-correlation} bound their covariance.
To restore a deleted bond, interpolate its strength. The derivative
of the corresponding mean is a conditional covariance with the bond
observable. Apply the conditional arbitrary-observable bound at radius
$\lfloor d/100\rfloor$ and dominate all weakened-system two-points by
the full conditional two-points. Averaging and Jensen again give the
same exponential scale. These are $L^1$ replacement bounds, which
suffice because the conditional means are bounded.
The neighborhood volumes and crossing-edge lists have polynomial
size, and small $d$ is covered by the trivial estimate. This proves
Equation~\eqref{eq:prelim-mixing}, with possibly enlarged $C,p$, because
$s\le X(b)$ by Equation~\eqref{eq:prelim-total-length}.

For completeness, the volume and partition-function conclusions use
the same uniformity under bond deletion. Cutting a collar around a
fixed support and integrating the coupling derivatives compares any
two large periodic volumes by a polynomial times the exponentially
small boundary error. The local periodic expectations are therefore
Cauchy, independently of the exhaustion, giving the periodic limit and
its uniqueness as in \Rref{pre:conclusion}.
For homogeneous rectangles, apply the common-cutout comparison of
\Rref{pre:part-12-3} once to horizontal bond expectations and once to
vertical bond expectations. The difference between the normalized
bond sums on the $n$-by-$w$ and $2n$-by-$2w$ tori is bounded by a
polynomial in $b,n,w$ times $e^{-c\min(n,w)/X(b)}$.
Integration of the coupling from zero to $b$ gives the upper bound in
Equation~\eqref{eq:prelim-doubling}; the comparison is uniform at every
intermediate coupling because all strengths remain in $[0,b]$.
Finally, the row and column transfer operators are nonnegative:
$e^{bq\cdot q'}$ is a positive kernel by its tensor-power expansion.
Twice applying $\Tr T^{2k}\le(\Tr T^k)^2$, in the two coordinate
directions, gives $Z_b(2n,2w)\le Z_b(n,w)^4$. This proves the lower
bound and completes Proposition~\ref{prop:preliminary}.

From Section~\ref{sec:free} onward, $L$ denotes one fixed sufficiently
large dyadic renormalization factor. Constants chosen before $L$ are
uniform for all sufficiently large $L$; constants denoted $C_L$ may
change after $L$ is fixed. The same convention applies after choosing
the terminal inverse coupling $H$. Fixed constants in positive-power
error estimates are understood in this order of choices.
